\subsection{The Case of Simple Rings}

Let A and B be simple rings and f a homomorphism from A to B. Let S be a simple A-module and T a simple B-module. We set

\[
i (f) = [ f ^ {*} (\mathrm{S}): \mathrm{T} ] = \mathrm{long} _ {\mathrm{B}} (f ^ {*} (\mathrm{S}));\tag{26}
\]

we call the cardinal \(i(f)\) the index of f.~When A is a subring of B and f is the canonical injection of A into B, we write \(i(\mathrm{B},\mathrm{A})\) instead of \(i(f)\) , and we call this cardinal the index of A in B. We define the height \(h(f)\) of f analogously:

\[
h (f) = [ f _ {*} (\mathrm{T}): \mathrm{S} ] = \mathrm{long} _ {\mathrm{A}} (f _ {*} (\mathrm{T})).\tag{27}
\]

When A is a subring of B and f the canonical injection of A into B, we write \(h(\mathrm{B}, \mathrm{A})\) for \(h(f)\) , and we call it the height of A in B.

The A-module S is monogenous, hence so is the B-module \(f^{*}(\mathrm{S}) = \mathrm{B} \otimes_{\mathrm{A}} \mathrm{S}\) . It follows that \(i(f)\) is finite and therefore that it is an integer. Let M be an A-module. Denote its length by a; then M is isomorphic to \(\mathrm{S}^{(\mathfrak{a})}\) . Consequently, the B-module \(f_{*}(\mathrm{M})\) is isomorphic to \(f_{*}(\mathrm{S})^{(\mathfrak{a})}\) . By the definition of \(i(f)\) , we therefore have

\[
\mathrm{long} _ {\mathrm{B}} (f ^ {*} (\mathrm{M})) = i (f) \mathrm{long} _ {\mathrm{A}} (\mathrm{M}).\tag{28}
\]

The Z-modules \(K_{0}(A)\) and \(K_{0}(B)\) are free of dimension 1, with respective bases {[}S{]} and {[}T{]}, and we have

\[
f ^ {*} ([ \mathrm{S} ]) = i (f) [ \mathrm{T} ].\tag{29}
\]

Let us take, in particular, \(M = A_{s}\) ; then \(f^{*}(A_{s}) = B \otimes_{A} A_{s}\) is isomorphic to \(B_{s}\) (II, §3, No.~4, p.~249), so

\[
i (f) = \mathrm{long(B)} / \mathrm{long(A)}.\tag{30}
\]

By Wedderburn's theorem (VIII, p.~120, Theorem 1), there exist integers \(m \geqslant 1\) and \(n \geqslant 1\) and fields D and E such that A is isomorphic to \(\mathbf{M}_m(\mathrm{D})\) and B to \(\mathbf{M}_n(\mathrm{E})\) . By formula (30), we have \(i(f) = n / m\) ; in particular, \(m\) divides \(n\) .

Let N be an B-module; denote its length by a. Then N is isomorphic to \(\mathrm{T}^{(\mathfrak{a})}\) , so the A-module \(f_{*}(\mathrm{N})\) is isomorphic to \(f_{*}(\mathrm{T})^{(\mathfrak{a})}\) ; by the definition of \(h(f)\) , we have

\[
\mathrm{long} _ {\mathrm{A}} (f _ {*} (\mathrm{N})) = h (f) \mathrm{long} _ {\mathrm{B}} (\mathrm{N}).\tag{31}
\]

We have seen (VIII, p.~124, Proposition 5) that f makes B into a free A-module and that all bases of this module have the same cardinal, denoted by \([B:A]_{s}\) and called the (left) degree of B over A. The A-module \(f_{*}(B_{s})\) is isomorphic to \(A_{s}^{[B:A]_{s}}\) , so has length \([B:A]_{s}\) long(A). By formula (30) and formula (31) applied to the specific case \(N=B_{s}\) , we therefore have

\[
[ \mathrm{B}: \mathrm{A} ] _ {s} = i (f) h (f).\tag{32}
\]

Now suppose that B is a finitely generated A-module, that is, that \([B:A]_{s}\) is finite. Then \(h(f)\) is finite by formula (32). We defined (VIII, p.~202) a group homomorphism \(f_{*}\) from \(K_{0}(B)\) to \(K_{0}(A)\) ; we have

\[
f _ {*} ([ \mathrm{T} ]) = h (f) [ \mathrm{S} ].\tag{33}
\]

Suppose that A and B are finite-dimensional algebras over a commutative field K and that f is K-linear. As before, there exist integers \(m \geqslant 1\) and \(n \geqslant 1\) , K-algebras D and E that are fields, and K-algebra isomorphisms from A to \(\mathbf{M}_{m}(\mathrm{D})\) and from B to \(\mathbf{M}_{n}(\mathrm{E})\) . Set \(d = [D : K]\) and \(e = [E : K]\) . We then have the relations

\[
[ \mathrm{A}: \mathrm{K} ] = m ^ {2} d, \qquad [ \mathrm{B}: \mathrm{K} ] = n ^ {2} e, \qquad [ \mathrm{B}: \mathrm{A} ] _ {s} = \frac {n ^ {2} e}{m ^ {2} d}
\]

and, by formulas (30) and (32), the relations

\[
i (f) = \frac {n}{m}, \qquad h (f) = \frac {n e}{m d}.
\]

When the field \(\mathrm{K}\) is algebraically closed, we have \(d = e = 1\) and therefore \(i(f) = h(f)\) and \([\mathrm{B}:\mathrm{A}]_s = i(f)^2\) .

Let A, B, and C be simple rings, and let \(f: A \to B\) and \(g: B \to C\) be homomorphisms. Let S be a simple A-module. The C-modules \((g \circ f)^{*}(S)\) and \(g^{*}(f^{*}(S))\) are isomorphic; therefore, by formulas (26) and (28), we have

\[
i (g \circ f) = \operatorname{long} _ {\mathrm{C}} \left(g ^ {*} \left(f ^ {*} (\mathrm{S})\right)\right) = i (g) \operatorname{long} _ {\mathrm{B}} \left(f ^ {*} (\mathrm{S})\right) = i (g) i (f).
\]

We can prove the equality \(h(g \circ f) = h(g)h(f)\) likewise. When A is a subring of B and B is a subring of C, and we take the canonical injections for f and g, these equalities give

\[
i (\mathrm{C}, \mathrm{A}) = i (\mathrm{C}, \mathrm{B}) i (\mathrm{B}, \mathrm{A}), \qquad h (\mathrm{C}, \mathrm{A}) = h (\mathrm{C}, \mathrm{B}) h (\mathrm{B}, \mathrm{A}).\tag{34}
\]

PROPOSITION 12. --- Let B be a simple ring and A a simple subring of B. Suppose that B is a finitely generated left A-module. Let M be a nonzero finitely generated left B-module. Set \(A' = \text{End}_{A}(M)\) and \(B' = \text{End}_{B}(M)\) . Then \(B'\) is a subring of \(A'\) , the rings \(A'\) and \(B'\) are simple, \(A'\) is a finitely generated left \(B'\) -module, and we have the equalities

\[
i (\mathrm{A} ^ {\prime}, \mathrm{B} ^ {\prime}) = h (\mathrm{B}, \mathrm{A}), \quad h (\mathrm{A} ^ {\prime}, \mathrm{B} ^ {\prime}) = i (\mathrm{B}, \mathrm{A}), \quad [ \mathrm{A} ^ {\prime}: \mathrm{B} ^ {\prime} ] _ {s} = [ \mathrm{B}: \mathrm{A} ] _ {s}.
\]

By Proposition 4 of VIII, p.~123, the ring \(\mathrm{A}'\) is simple, \(\mathrm{M}\) is an \(\mathrm{A}'\) -module of finite length, and we have

\[
\mathrm{long} _ {\mathrm{A}} (\mathrm{M}) = \mathrm{long} (\mathrm{A} ^ {\prime}), \quad \mathrm{long} _ {\mathrm{A} ^ {\prime}} (\mathrm{M}) = \mathrm{long} (\mathrm{A}).
\]

For the same reasons, the ring \(B'\) is simple, and we have

\[
\operatorname{long} _ {\mathrm{B}} (\mathrm{M}) = \operatorname{long} \left(\mathrm{B} ^ {\prime}\right), \quad \operatorname{long} _ {\mathrm{B} ^ {\prime}} (\mathrm{M}) = \operatorname{long} (\mathrm{B}).
\]

By formulas (31) and (30), we therefore have

\[
h (\mathrm{B}, \mathrm{A}) = \operatorname{long} _ {\mathrm{A}} (\mathrm{M}) / \operatorname{long} _ {\mathrm{B}} (\mathrm{M}) = \operatorname{long} \left(\mathrm{A} ^ {\prime}\right) / \operatorname{long} \left(\mathrm{B} ^ {\prime}\right) = i \left(\mathrm{A} ^ {\prime}, \mathrm{B} ^ {\prime}\right);
\]

the equality \(h(\mathrm{A}^{\prime},\mathrm{B}^{\prime}) = i(\mathrm{B},\mathrm{A})\) can be established analogously. From this, we deduce

\[
[ \mathrm{A} ^ {\prime}: \mathrm{B} ^ {\prime} ] _ {s} = i (\mathrm{A} ^ {\prime}, \mathrm{B} ^ {\prime}) h (\mathrm{A} ^ {\prime}, \mathrm{B} ^ {\prime}) = h (\mathrm{B}, \mathrm{A}) i (\mathrm{B}, \mathrm{A}) = [ \mathrm{B}: \mathrm{A} ] _ {s}
\]

using formula (32). In particular, \(A'\) is a finitely generated left \(B'\) -module.

